\documentclass[12pt]{article}
\usepackage[T1]{fontenc}
\usepackage[utf8]{inputenc}
\usepackage{lmodern}
\usepackage{textcomp}
\usepackage{enumitem}
\usepackage{geometry}
\geometry{margin=1in}
\begin{document}
\begin{center}
Answer Key - Constitutional Law - Undergraduate Level Assessment 21 \
\small (Answers are ordered exactly as in the question paper.)
\end{center}

\section*{Multiple Choice}
\begin{enumerate}[label=\arabic*.]
\item \textbf{Answer:} B
\\ \textit{Options:} A) interpret; B) make; C) enforce; D) override
\\ \textit{Note:} The legislative branch of the United States government, comprised of Congress, is explicitly granted the power to make laws as outlined in Article I of the Constitution.
\item \textbf{Answer:} D
\\ \textit{Options:} A) That natural law corresponds to positive law.; B) That Hume is a legal positivist.; C) That syllogistic logic is false.; D) That human reason can help us to determine what constitutes a worthwhile life.
\\ \textit{Note:} John Finnis argues that human reason is capable of discerning the fundamental goods that contribute to a worthwhile life, thereby rejecting Hume's view that practical reason is merely a tool for fulfilling desires without a rational basis for determining what is ultimately valuable.
\item \textbf{Answer:} C
\\ \textit{Options:} A) There is no distinction between the two forms of legal reasoning.; B) Judges are appointed to interpret the law, not to make it.; C) It is misleading to pigeon-hole judges in this way.; D) Judicial reasoning is always formal.
\\ \textit{Note:} The criticism that it is misleading to pigeon-hole judges suggests that legal reasoning is often more complex and nuanced than a simple dichotomy allows, failing to account for the diverse approaches judges may take in different contexts.
\item \textbf{Answer:} D
\\ \textit{Options:} A) Conflict.; B) Love.; C) War.; D) Reciprocity.
\\ \textit{Note:} The correct answer is reciprocity because Malinowski emphasized that the social structure of the Trobriand Islanders was fundamentally based on the mutual exchange of goods and services, which reinforced social bonds and community cohesion.
\item \textbf{Answer:} B
\\ \textit{Options:} A) John Rawls; B) Stammler; C) Jerome Hall; D) John Finns
\\ \textit{Note:} Stammler is recognized for advocating a theory of natural law that emphasizes its adaptability to changing societal conditions and values, thus embodying the concept of "natural law with a variable content."
\end{enumerate}

\section*{True / False}
\begin{enumerate}[label=\arabic*.]
\setcounter{enumi}{5}
\item \textbf{Answer:} False
\\ \textit{Options:} A) True; B) False
\\ \textit{Note:} Kant emphasized the intrinsic value of the individual and personal autonomy, asserting that moral actions must arise from rational will and universal moral laws rather than solely from societal norms.
\item \textbf{Answer:} False
\\ \textit{Options:} A) True; B) False
\\ \textit{Note:} The statement is false because the revival of natural law principles in the 20 th century was primarily influenced by post-World War II developments and philosophical movements, rather than the economic downturn of the 1930 s, which instead led to an emphasis on economic regulation and legal positivism.
\item \textbf{Answer:} True
\\ \textit{Options:} A) True; B) False
\\ \textit{Note:} The statement is true because the Commerce Clause of the U.S. Constitution grants Congress the authority to regulate commerce with foreign nations and among the states, thereby enabling it to enact federal laws governing these areas.
\item \textbf{Answer:} True
\\ \textit{Options:} A) True; B) False
\\ \textit{Note:} In Hohfeld's scheme, power allows one party to create or alter legal obligations, while liability represents a legal obligation that one party has towards another, making them fundamentally incompatible as one cannot simultaneously hold the power to impose a duty while being subject to that duty.
\item \textbf{Answer:} False
\\ \textit{Options:} A) True; B) False
\\ \textit{Note:} The statement is false because not all acts performed by a state are considered jure imperii; specifically, commercial activities or acts outside the exercise of sovereign authority may not attract immunity.
\end{enumerate}

\section*{Long Form Response}
\begin{enumerate}[label=\arabic*.]
\setcounter{enumi}{10}
\item \textbf{Answer:} The primary focus of the European Convention on Human Rights (ECHR) is to safeguard civil and political rights across its member states, ensuring that individuals can exercise fundamental freedoms such as the right to life, freedom from torture, and the right to a fair trial. These rights are enshrined in the Convention's articles, which have significant implications for constitutional law, as they create binding obligations for states to uphold these rights and provide effective remedies for violations. This legal framework not only enhances the protection of individual liberties but also influences national constitutions and legal systems, promoting a culture of human rights within Europe. Consequently, the ECHR serves as a crucial instrument in advancing democratic principles and accountability in governance, thereby reinforcing the rule of law across the continent.
\\ \textit{Note:} correct because it accurately identifies the ECHR's role in protecting civil and political rights through its binding articles, which have a profound impact on the constitutional obligations of member states to ensure individual freedoms and provide legal remedies.
\item \textbf{Answer:} Emile Durkheim's concepts of mechanical and organic solidarity describe how societies maintain cohesion and social order. Mechanical solidarity, characteristic of traditional societies, arises from the homogeneity of individuals and leads to repressive laws, which serve to enforce conformity and punish deviance. In contrast, organic solidarity, typical of modern, complex societies, is based on the interdependence of individuals performing specialized roles, resulting in restitutive laws that focus on restoring social equilibrium rather than punishing offenders. Thus, the nature of law in a society reflects its underlying solidarity; repressive laws maintain uniformity in mechanically solidary societies, while restitutive laws facilitate cooperation and adjustment in organically solidary societies.
\\ \textit{Note:} The answer accurately describes Durkheim's distinction between mechanical and organic solidarity, linking mechanical solidarity to repressive laws that enforce conformity and organic solidarity to restitutive laws that reflect the interdependence of specialized roles in modern societies.
\end{enumerate}

\vfill
\noindent\textit{End of Answer Key}
\end{document}